\pdfminorversion=4
\pdfobjcompresslevel=0
\documentclass[ignorenonframetext,]{beamer}
\setbeamertemplate{caption}[numbered]
\setbeamertemplate{caption label separator}{:}
\setbeamercolor{caption name}{fg=normal text.fg}
\usepackage{amssymb,amsmath}
\usepackage{colortbl}
\usepackage{booktabs}
\usepackage{pifont}
\usepackage{graphicx}
\usepackage{ifxetex,ifluatex}
\usepackage{fixltx2e} % provides \textsubscript
\usepackage{lmodern}
\usepackage{pdfpages}
%\setbeamercolor{background canvas}{bg=}
\makeatletter\beamer@ignorenonframefalse\makeatother
% see https://tex.stackexchange.com/questions/96581/incorrect-frame-numbers-when-pages-included-using-pdfpages-in-beamer
\def\getpdfpages#1#2{\begingroup
	\setbeamercolor{background canvas}{bg=}
	\includepdf[pages={#1}, trim=0mm 1mm 0mm 0mm,%
	pagecommand={\global\setcounter{framenumber}{%
			\the\numexpr\value{framenumber}+1\relax}%
		\expandafter\def\expandafter\insertshorttitle\expandafter{%
			\insertshorttitle\hfill\insertframenumber\,/\,\inserttotalframenumber}}%
	]{#2}
	\endgroup}

\ifxetex
  \usepackage{fontspec,xltxtra,xunicode}
  \defaultfontfeatures{Mapping=tex-text,Scale=MatchLowercase}
  \newcommand{\euro}{€}
\else
  \ifluatex
    \usepackage{fontspec}
    \defaultfontfeatures{Mapping=tex-text,Scale=MatchLowercase}
    \newcommand{\euro}{€}
  \else
    \usepackage[T1]{fontenc}
    \usepackage[utf8]{inputenc}
      \fi
\fi
% use upquote if available, for straight quotes in verbatim environments
\IfFileExists{upquote.sty}{\usepackage{upquote}}{}
% use microtype if available
\IfFileExists{microtype.sty}{\usepackage{microtype}}{}
\usepackage{url}

\providecommand{\tightlist}{%
  \setlength{\itemsep}{0pt}\setlength{\parskip}{0pt}}

% Comment these out if you don't want a slide with just the
% part/section/subsection/subsubsection title:
\AtBeginPart{
  \let\insertpartnumber\relax
  \let\partname\relax
  \frame{\partpage}
}
\AtBeginSection{
  \let\insertsectionnumber\relax
  \let\sectionname\relax
  \frame{\sectionpage}
}
\AtBeginSubsection{
  \let\insertsubsectionnumber\relax
  \let\subsectionname\relax
  \frame{\subsectionpage}
}

\setlength{\parindent}{0pt}
\setlength{\parskip}{6pt plus 2pt minus 1pt}
\setlength{\emergencystretch}{3em}  % prevent overfull lines
\setcounter{secnumdepth}{0}
\usepackage[english]{babel}
\usepackage{dsfont}
\usepackage{verbatim}
\usepackage{amsmath}
\usepackage{amssymb}
\usepackage{amsthm}
\usepackage{amsfonts}
\usepackage{csquotes}
\usepackage{multirow}
\usepackage{longtable}
\usepackage{enumerate}
% \usepackage[absolute,overlay]{textpos}
\usepackage{psfrag}
\usepackage{algorithm}
\usepackage{algpseudocode}
\usepackage{eqnarray}
\usepackage{arydshln}
\usepackage{tabularx}
\usepackage{placeins}
\usepackage{tikz}
% see https://tex.stackexchange.com/questions/32918/how-can-i-decrease-spaces-between-equations
%\usepackage{setspace}
\usepackage[nodisplayskipstretch]{setspace}
\setstretch{1.0}
\usetikzlibrary{shapes,arrows,automata,positioning,calc}
\usepackage{subfig}
% \usepackage{paralist}
\usepackage{graphicx}
\usepackage{array}
\usepackage{framed}
\usepackage{excludeonly}
\usepackage{fancyvrb}
\usecolortheme{dove}
% \usefonttheme{serif}
\usepackage{xfrac}
\usepackage{xcolor}
\usepackage{mdframed}
\usepackage{caption}
\captionsetup[figure]{labelformat=empty}
\usepackage{transparent}
\usepackage{bm}

% URL color:
\definecolor{myorange}{RGB}{80,224,176}
% Theme color of metropolis
\definecolor{metropolis_theme_color}{RGB}{38,38,38}
% Upper and lower rules of code chunks
\definecolor{darkblue}{RGB}{38,38,38}
% Font color of footer:
\definecolor{mygrey}{RGB}{240,240,240}

\definecolor{accents}{RGB}{80,224,176}

\usepackage{hyperref}
\hypersetup{
    colorlinks = true,
    linkcolor = {black},
    urlcolor = {accents},
    linkbordercolor = {white}
}


%% %%%%%%%%%%%%%%%%%%%%%%%%%%%%%%%%%%%%%%%%%%%%%%%%%%%%%%%%%%%%%%%
%% METROPOLIS THEME + CUSTOMIZATIONS
%% %%%%%%%%%%%%%%%%%%%%%%%%%%%%%%%%%%%%%%%%%%%%%%%%%%%%%%%%%%%%%%%

%\usetheme{metropolis}
%\usetheme{metropolis}
%\useoutertheme{metropolis}
%\useinnertheme{metropolis}
\setbeamercolor{background canvas}{bg=white}
\AtBeginSection{
  \let\insertsectionnumber\relax
  \let\sectionname\relax
    \begin{frame}[c,plain,noframenumbering]
    \sectionpage
    \small{\tableofcontents[currentsection]}
    \end{frame}
}

% remove space of center enviroment which often occurs around includegraphics produced from knitr plots
% https://tex.stackexchange.com/questions/74212/how-can-i-change-the-whitespace-above-and-below-center
\let\oldcenter\center
\let\oldendcenter\endcenter
\renewenvironment{center}{\setlength\topsep{0pt}\oldcenter}{\oldendcenter}

% reduce margins / border
% \setbeamersize{text margin left=10pt,text margin right=10pt}

% Reduce space between code chunks and code output in rmarkdown beamer presentation: https://stackoverflow.com/questions/35734525/reduce-space-between-code-chunks-and-code-output-in-rmarkdown-beamer-presentatio
%\setlength{\parskip}{0pt}
%\setlength{\OuterFrameSep}{-20pt}
\makeatletter
\preto{\@verbatim}{\parskip=0pt \topsep=0pt \partopsep=0pt}
\makeatother

%% Add LMU template:
%% ---------------------------------------------------------------
%\usepackage{./../../_setup/lmu-lecture}
\usepackage{../../style/lmu-lecture}
% \usefonttheme{structurebold}
% \RequirePackage[T1]{fontenc}
% \RequirePackage[scaled=0.92]{helvet}   %% Helvetica for sans serif
% \defbeamertemplate*{frametitle}{lecture}[1][left]
% {
%   \ifbeamercolorempty[bg]{frametitle}{}{\nointerlineskip}%
%   \@tempdima=\textwidth%
%   \advance\@tempdima by\beamer@leftmargin%
%   \advance\@tempdima by\beamer@rightmargin%
%   \begin{beamercolorbox}[sep=0.2cm,#1,wd=\the\@tempdima]{frametitle}
%     \if@tempswa\else\csname beamer@fte#1\endcsname\fi%
%     {\usebeamerfont{frametitle}\rule[-0.5ex]{0pt}{2.3ex}\insertframetitle\par}%
%     \if@tempswa\else\vskip-.2cm\fi% set inside beamercolorbox... evil here...
%   \end{beamercolorbox}%
% }
% \def\beamer@fteright{\vskip0.35cm\advance\leftskip by 1.7cm\advance\rightskip by1.7cm}
% \setbeamertemplate{frametitle continuation}[from second][{\small/~\insertcontinuationcount}]
% % ------------------------------------------------------------------------
% % Geometry
% \setbeamersize{text margin left=0.8cm,text margin right=0.8cm}
% \pgfdeclarehorizontalshading{footlineshade}{4mm}{%
%   color(0pt)=(black);%
%   color(1.0\paperwidth)=(structure!50!black)}
% % redefine \ref (it has been redefined somewhere by the beamerclass)
% \def\lectureref#1{\expandafter\@setref\csname r@#1\endcsname\@firstoftwo{#1}}
% % counter for framenumber for current lecture
% \newcounter{lectureframenumber}
% % adjust framenumbers for lecture (check whether reference is already defined)
% \def\addjustlectureframenumber#1{\ifx#1\relax\else%
%   \addtocounter{lectureframenumber}{-\lectureref{lect:@@\thelecture}}\fi}

% % ------------------------------------------------------------------------
% % Navigation symbols
% \setbeamertemplate{navigation symbols}{}
% % ------------------------------------------------------------------------
% % No head lines
% \defbeamertemplate*{headline}{lecture theme}{}
% \setbeamertemplate{frametitle}{\expandafter\uppercase\expandafter\insertframetitle}








%% Add footer:
%% ---------------------------------------------------------------
\setbeamertemplate{footline}[text line]{%
    \noindent\hspace*{\dimexpr-\oddsidemargin-1in\relax}%
     \colorbox{white}{
     \makebox[\dimexpr\paperwidth-2\fboxsep\relax]{
     \color{black}
     \begin{minipage}[c][4.5ex][c]{0.5\linewidth}
       \secname
     \end{minipage}
     \hfill\begin{minipage}[c][4.5ex][c]{0.5\linewidth}
       \flushright
       \insertframenumber{}~/~\inserttotalframenumber~~
     \end{minipage}
     }}%
  \hspace*{-\paperwidth}
}

%% Add titlegraphic:
%% ---------------------------------------------------------------
% \titlegraphic{
%   \transparent{0.7}
%   \includegraphics[width=5cm]{./../../_setup/images/logos/eds_logo_black.png}
% }

%% %%%%%%%%%%%%%%%%%%%%%%%%%%%%%%%%%%%%%%%%%%%%%%%%%%%%%%%%%%%%%%%
%% NICER CODE BACKGROUND
%% %%%%%%%%%%%%%%%%%%%%%%%%%%%%%%%%%%%%%%%%%%%%%%%%%%%%%%%%%%%%%%%

% Define Shaded if not defined already through knitr:
%\definecolor{shadecolor}{RGB}{228,228,228}
\makeatletter
\@ifundefined{Shaded}{%
\newenvironment{Shaded}{\begin{snugshade}}{\end{snugshade}}%
}{}
\makeatother

\renewenvironment{Shaded}{
  %\vspace{5pt}
  \begin{mdframed}[backgroundcolor=black!8,
                   linecolor=darkblue,
                   rightline=false,
				   leftline=false]}{
  \end{mdframed}}


%% %%%%%%%%%%%%%%%%%%%%%%%%%%%%%%%%%%%%%%%%%%%%%%%%%%%%%%%%%%%%%%%
%% CHANGE URL COLOR
%% %%%%%%%%%%%%%%%%%%%%%%%%%%%%%%%%%%%%%%%%%%%%%%%%%%%%%%%%%%%%%%%

\renewcommand\UrlFont{\color{accents}}
% \hypersetup{
%     colorlinks=true,
%     linkcolor=myorange
% }

%% %%%%%%%%%%%%%%%%%%%%%%%%%%%%%%%%%%%%%%%%%%%%%%%%%%%%%%%%%%%%%%%
%% INCLUDE LATEX MATH + OTHER CUSTOMIZATIONS
%% %%%%%%%%%%%%%%%%%%%%%%%%%%%%%%%%%%%%%%%%%%%%%%%%%%%%%%%%%%%%%%%


% Include latex-math
\DeclareOldFontCommand{\sf}{\normalfont\sffamily}{\mathsf}

% \input{./../../latex-math/basic-math.tex}
% \input{./../../latex-math/basic-ml.tex}
% \input{./../../latex-math/ml-bagging.tex}
% \input{./../../latex-math/ml-boosting.tex}
% \input{./../../latex-math/ml-gp.tex}
% \input{./../../latex-math/ml-mbo.tex}
% \input{./../../latex-math/ml-nn.tex}
% \input{./../../latex-math/ml-svm.tex}
% \input{./../../latex-math/ml-trees.tex}
% \InputIfFileExists{./../../latex-math/ml-lm.tex}{}{}
% \InputIfFileExists{./../../latex-math/ml-automl.tex}{}{}
% Additional commands:
\newcommand{\abs}[1]{\ensuremath{\left| #1 \right|}}
%\renewcommand{\pikN}{\hat{\pi}^\Np} % This collides with z



% basic latex stuff
\newcommand{\pkg}[1]{{\fontseries{b}\selectfont #1}} %fontstyle for R packages
\newcommand{\lz}{\vspace{0.5cm}} %vertical space

\newcommand{\oneliner}[1] % Oneliner for important statements
{\begin{block}{}\begin{center}\begin{Large}#1\end{Large}\end{center}\end{block}}

\let\code=\texttt
\let\proglang=\textsf

\setkeys{Gin}{width=0.9\textwidth}

%% %%%%%%%%%%%%%%%%%%%%%%%%%%%%%%%%%%%%%%%%%%%%%%%%%%%%%%%%%%%%%%%
%% EXERCISE COUNTER
%% %%%%%%%%%%%%%%%%%%%%%%%%%%%%%%%%%%%%%%%%%%%%%%%%%%%%%%%%%%%%%%%

% Command for an automatic exercise counter for the exercise sheets:
\newcounter{Exercise}
\newcommand{\exercise}{
  \stepcounter{Exercise}
  \ifnum \theExercise = 1
    \Large{\textbf{Exercise \theExercise}} \\
  \else
    \vspace{\baselineskip} \Large{\textbf{Exercise \theExercise}} \\
  \fi
}

% Command for an automatic exercise counter for the courses:
\newcounter{CourseExercise}
\newcommand{\oneExercise}{
  \stepcounter{CourseExercise}
    \begin{center}
      \huge{Exercise \Roman{CourseExercise}}
    \end{center}
}
\newcommand{\twoExercises}{
  \stepcounter{CourseExercise}
    \begin{center}
      \huge{Exercise \Roman{CourseExercise} \& \stepcounter{CourseExercise} \Roman{CourseExercise}}
    \end{center}
}
\newcommand{\threeExercises}{
  \stepcounter{CourseExercise}
    \begin{center}
      \huge{Exercise \Roman{CourseExercise}, \stepcounter{CourseExercise} \Roman{CourseExercise} \& \stepcounter{CourseExercise} \Roman{CourseExercise}}
    \end{center}
}

\newcommand{\resetExercises}{\setcounter{CourseExercise}{0}}

\usetikzlibrary{shapes,arrows,automata,positioning,calc,chains,trees, shadows}
\tikzset{
  %Define standard arrow tip
  >=stealth',
  %Define style for boxes
  punkt/.style={
    rectangle,
    rounded corners,
    draw=black, very thick,
    text width=6.5em,
    minimum height=2em,
    text centered},
  % Define arrow style
  pil/.style={
    ->,
    thick,
    shorten <=2pt,
    shorten >=2pt,}
}
\usepackage{subfig}
%\beamerdefaultoverlayspecification{<+->}

\title{Data Science Lifecycle}
\date{}

\begin{document}
\begin{frame}[plain]
\titlepage
\end{frame}


\begin{frame}[plain]{}
\phantomsection\label{section}
\end{frame}

\section{The Data Science Lifecycle}\label{the-data-science-lifecycle}

\begin{frame}{Learning Objective}
    \begin{itemize}
        \item overview over data science lifecycle
        \item become familiar with concepts and know what to look up in case you need it
        \item have some structure when it comes to completing a data science task
    \end{itemize}
\end{frame}

\begin{frame}{The Data Science Lifecycle}
\phantomsection\label{the-data-science-lifecycle-1}
\textbf{Lifecycle Management}: Process of managing a product lifecycle from inception, through engineering, design and manufacturing to deployment and eventual disposal.

%https://docs.google.com/presentation/d/1qiK0ShkIEbLwAzhEME_Yu7hz3lhnZhl2/edit?slide=id.p5#slide=id.p5
\begin{figure}[h]
\includegraphics[width=8cm]{slides/01_intro/DSLifecycle1.png}
\end{figure}

\end{frame}

\begin{frame}{The Data Science Lifecycle}
\phantomsection\label{the-data-science-lifecycle-1}

%https://docs.google.com/presentation/d/1qiK0ShkIEbLwAzhEME_Yu7hz3lhnZhl2/edit?slide=id.p5#slide=id.p5
\begin{figure}[h]
\includegraphics[width=11cm]{slides/01_intro/DSLifecycle2.png}
\end{figure}

\end{frame}


\section{Business Understanding}\label{business-understanding-1}

\begin{frame}{Overview Business Understanding}
%https://docs.google.com/presentation/d/1qiK0ShkIEbLwAzhEME_Yu7hz3lhnZhl2/edit?slide=id.p5#slide=id.p5
\begin{figure}[h]
\includegraphics[width=9cm]{slides/01_intro/BusinessUnderstandingOverview.png}
\end{figure}
\end{frame}

\begin{frame}{Requirements Analysis and Constraints}
\phantomsection\label{requirements-analysis}
\textbf{Key question}: What do you want to achieve in the end?
\begin{itemize}
\tightlist
\item
  Define the \textbf{result} of your work:
  \begin{itemize}
  \tightlist
  \item
    Simple report document?
  \item
    Web-app / dashboard?
  \item
    Machine Learning pipelines for internal/productive use?\\
    How does the model need to scale?
  \end{itemize}
\item Identify constraints (--> Summarized in \textbf{Service Level Agreements})
    \begin{itemize}
      \tightlist
      \item
        Ethical questions: We do not want to use the client's address
      \item
        Legal constraints: We are not allowed to use a job applicant's
        gender as a feature to automatically reject the application
      \item
        Reproducibility: The analysis must be fully reproducible in one step
      \item
        Explainability: We must explain all model decisions to the user
        \item Prediction Latency: We need a prediction in real-time
        \item Model Size: The model may only occupy a certain memory space
      \end{itemize}
\end{itemize}
\end{frame}

\begin{frame}{Non-ML Baseline}
\phantomsection\label{non-ml-baseline}
\textbf{Key question}: Do you really need ML?
\begin{itemize}
\tightlist
\item
  You don't have to use ML just because it's hyped.
\item
  What does the current solution look like?
\item
  Are other solutions possible that don't need ML?
\end{itemize}
\end{frame}

\begin{frame}{Choosing the Right Problem Frame}
\phantomsection\label{choosing-the-right-problem-frame}
\textbf{Key question}: Is it classification, regression or something else?
\begin{itemize}
\tightlist
\item
  A highly imbalanced problem: Frame it as classification, or anomaly
  detection?
\item
  Predictive maintenance: Frame it as classification (probability of
  failure), or regression (expected remaining lifetime)?
\item
  There answer is not clear, but depends on your application. It's
  important to think about these problems!
\end{itemize}
ML-tasks you can select from:
%https://docs.google.com/presentation/d/1qiK0ShkIEbLwAzhEME_Yu7hz3lhnZhl2/edit?slide=id.p5#slide=id.p5
\begin{figure}[h]
\includegraphics[width=7cm]{slides/01_intro/MLTasks.png}
\end{figure}
\end{frame}


\begin{frame}{KPIs and Business Objectives}
\phantomsection\label{translating-kpis-into-performance-measures}
\textbf{KPI}: Key Performance Indicator

\textbf{Key question}: How does the business side measure success/ performance of the ML solution? How can we translate this to a ML performance measure?

\begin{itemize}
\tightlist
\item
  Spam detection:

\begin{itemize}
  \tightlist
  \item
    \emph{Not} OK to misclassify a job application as spam (positive
    class)
  \item
    A spam mail in your inbox once in a while is no problem\\
    \(\rightarrow\) \emph{minimize the false positive rate},
    i.e.~maximize precision
  \end{itemize}
\item
  Airport security:

  \begin{itemize}
  \tightlist
  \item
    \emph{Not} OK to misclassify a weapon (positive class) as laptop
  \item
    A false alarm once in a while is no problem\\
    \(\rightarrow\) \emph{minimize the false negative rate},
    i.e.~maximize recall
  \end{itemize}
\end{itemize}

\end{frame}

\begin{frame}{Multiple objectives}
\phantomsection\label{multiple-objectives}
\textbf{Key question}: How to optimize multiple objectives?
\begin{itemize}
\tightlist
\item
  What if a choice increases objective 1, but decreases objective 2? Is
  it better?
\item i.e. in some case maximizing AUC leads to more model complexity and less explainability
\item
  There are two common ways to deal with multiple objectives:

  \begin{itemize}
  \tightlist
  \item
    \textbf{Scalarization} in a \emph{compound measure}:
    \(obj_1 * obj_2\)
  \item
    \textbf{Multi-objective Analysis} %TODO add link for further reading
  \end{itemize}
\end{itemize}
\end{frame}



\begin{frame}{Performance Measures}
\phantomsection\label{performance-measures}

\textbf{Key question}: what performance measure suits best for your solution?

\begin{itemize}
\tightlist
    \item \textbf{Precision and Recall}%: $Precision = \frac{TP}{TP+FP}$, $Recall = \frac{TP}{TP+FN}$
    \item to capture the tradeoff \textbf{F1-score} $F1 = \frac{2*Recall*Precision}{Recall + Precision}$
    \item \textbf{Accuracy}: %$Acc = \frac{TP+TN}{TP+TN+FP+FN}$
    for highly imbalanced data, the accuracy
  (percentage of correct predictions) is not a good measure
  \item I.e. in an HIV test, where 99.9\% of the population is healthy, a senseless
  test that predicts ``healthy'' \emph{all the time}, automatically has
  an accuracy of 99.9\%
\end{itemize}

%https://docs.google.com/presentation/d/1qiK0ShkIEbLwAzhEME_Yu7hz3lhnZhl2/edit?slide=id.p5#slide=id.p5
\begin{figure}[h]
\includegraphics[width=10cm]{slides/01_intro/ConfusionMatrix.png}
\end{figure}
\end{frame}




\begin{frame}{Validation Strategy}
\phantomsection\label{validation-strategy}

%TODO include graphics of train-test-tune-validate process
%TODO integrate all other validation points from later in the slides here and don"t come back

\begin{itemize}
\tightlist
\item
  How to validate your model? Understand what you want to generalize to

  \begin{itemize}
  \tightlist
  \item
    blocking factors
  \item
    temporal effects
  \item
    spatial effects
  \end{itemize}
\item
  Validate it just like you would use it in real life:

  \begin{itemize}
  \tightlist
  \item
    A time series model: You'd build it in 2020 with data from
    2016--2020, then use it in 2021
  \item
    Validate it with the same strategy: Train on 2016--2019, validate on
    2020
  \end{itemize}
\end{itemize}
\end{frame}





\section{Data Acquisition and
Understanding}\label{data-acquisition-and-understanding-1}

\begin{frame}{Data Source}
\phantomsection\label{data-source}
\begin{itemize}
\tightlist
\item
  In an ideal world, you would have only one data source, such as an
  enterprise data warehouse
\item
  Realistically, most businesses combine multiple internal datasets as
  well as a number of external sources
\end{itemize}
\end{frame}

\begin{frame}{Data Source - Internal vs.~external data sources}
\phantomsection\label{data-source---internal-vs.-external-data-sources}
\begin{itemize}
\tightlist
\item
  \textbf{Internal data sources} come from inside your organization and
  often contain data relating to your company (operations, finance,
  personnel).
\item
  \textbf{External data} refers to any type of data that has been
  captured, processed, and provided from outside the company.

  \begin{itemize}
  \tightlist
  \item
    External data includes data such as weather forecasts by an external
    provider's API, or a HTML website you scrape periodically.
  \end{itemize}
\item
  Internal data is often preferred because you own it and you are
  responsible for the format and continuous availability of the data.
\item
  External data providers can change the format or discontinue their
  service.
\end{itemize}
\end{frame}

\begin{frame}{Data Source - Structured vs.~unstructured data}
\phantomsection\label{data-source---structured-vs.-unstructured-data}
\begin{itemize}
\tightlist
\item
  Structured data means the traditional data set as you know it from an
  Excel table: Dates, Names, numbers, and strings.

  \begin{itemize}
  \tightlist
  \item
    It can easily be stored in a database, searched, and analyzed
  \end{itemize}
\item
  Unstructured data includes things like text documents and images. They
  are not straightforward to work with, but recent developments allowed
  using them in more and more powerful ways.
\end{itemize}
\end{frame}

\begin{frame}{Data Source - Database vs.~Files}
\phantomsection\label{data-source---database-vs.-files}
\begin{itemize}
\tightlist
\item
  Your data may come from many sources, and in various formats
  %TODO think of startegies to harmonize
  %TODO think of possible changes and how your model could adapt
\item
  \textbf{Text files} are easy to share and can be read by a human

  \begin{itemize}
  \tightlist
  \item
    However, column types are not specified or enforced. Excel will just
    assume your column is a date.
  \end{itemize}
\item
  \textbf{Binary files} are usually smaller in size, and additional
  information can be stored in them

  \begin{itemize}
  \tightlist
  \item
    However, they are often harder to transfer between applications or
    colleagues
  \end{itemize}
\item
  \textbf{Excel files} combine the disadvantages of text files and
  binary files. Avoid them.
\item
  \textbf{Databases} combine most of the advantages discussed: The
  schemas are clearly defined and most software can easily grab and
  process data from them.
\end{itemize}
\end{frame}

\begin{frame}{Data Source - On-premise vs.~cloud}
\phantomsection\label{data-source---on-premise-vs.-cloud}
\begin{itemize}
\tightlist
\item
  Determine whether your data is best stored locally in your company's
  own data center, or in a cloud environment such as Amazon Web Services
\item
  The decision will be influenced by data protection laws, costs
  (upfront, maintenance/upgrading, and recurring), and experience of
  your engineering team
\end{itemize}
\end{frame}

\begin{frame}{Data Access at Deployment Time}
\phantomsection\label{data-access-at-deployment-time}
\begin{itemize}
\tightlist
\item
  How to feed the \emph{new} data to your model when it's in use?
\item
  This is often a separate workflow than your training data pipeline
\item
  Plan ahead: Think about these questions in the beginning
\end{itemize}
\end{frame}

\begin{frame}{Data Exploration}
\phantomsection\label{data-exploration}
\begin{itemize}
\tightlist
\item
  First, understand your data in terms of their features
\item
  You should obtain or create a \emph{data dictionary}. For each
  feature:

  \begin{itemize}
  \tightlist
  \item
    What does the feature represent?
  \item
    In what unit is it measured? What's the possible range of values?
  \item
    For factors: What levels are possible?
  \item
    How are missing values encoded?
  \end{itemize}
\item
  Compute summary statistics for each variable
\item
  Plot a few uni- and bivariate diagrams
\item
  This way, you can find outliers or implausible values, and get a first
  feeling for imbalances, correlations, etc.
\end{itemize}
\end{frame}

\begin{frame}{Data Quality}
\phantomsection\label{data-quality}
\begin{itemize}
\tightlist
\item
  The quality of your data is crucial and will impact the model and your
  solution.

  \begin{itemize}
  \tightlist
  \item
    \emph{Garbage in, garbage out}
  \end{itemize}
\item
  How many missing values are in your columns?
\item
  Are there incorrect or implausible values? (You need domain knowledge
  for this)
\item
  Parse date/time values into an appropriate format
\item
  Remove nearly constant features
\item
  Remove redundant and obsolete features
\end{itemize}
\end{frame}

%TODO add these headlines for data quality (completeness, duplicates, consistency, accuracy, and conformity)

\begin{frame}{Datasheets}
\phantomsection\label{datasheets}
\begin{itemize}
\tightlist
\item
  In the electronics industry, every component comes with a
  \emph{datasheet} that describes everything from key numbers to
  recommended usage
\item
  Ideally, you document your dataset with an accompanying datasheet in a
  similar manner
\item
  See also: \href{https://arxiv.org/pdf/1803.09010.pdf}{Datasheets for
  Datasets}
\end{itemize}
\end{frame}

\begin{frame}[fragile]{Data/Label Leakage}
\phantomsection\label{datalabel-leakage}
%TODO include split strategies (graphics)
%TODO harmonize with validation?
\begin{itemize}
\tightlist
\item
  Data Leakage (or Label Leakage) describes an unwanted situation where
  you use data for training that you won't have available at prediction
  time.
\item
  Simple example (which \emph{does} happen in reality):

  \begin{itemize}
  \tightlist
  \item
    Your data has a feature \texttt{row\_number}, and the first 300 rows
    have the target class 0, the last 100 rows have the target class 1.
  \end{itemize}
\item
  Trickier examples:

  \begin{itemize}
  \tightlist
  \item
    Transforming your data (PCA, min/max, n-grams, etc.) with the
    \emph{entire} data set instead of just the training data
  \item
    If your data contains duplicate rows, having one row in the training
    data and one in the test data will lead to overconfidence even on
    the test performance
  \end{itemize}
\item
  It's tempting because data leakage usually improves the training
  performance of your model
\end{itemize}
\end{frame}

\begin{frame}{Data Quality / Biased Data}
\phantomsection\label{data-quality-biased-data}
\begin{itemize}
\tightlist
\item
  You need to ensure your data corresponds to the data generating
  process
\item
  Simple example: A ``random'' political survey, if held in the city
  center on a Tuesday at 10am, will overrepresent the opinion of senior
  citizens and students, and underrepresent people currently at work
\item
  \textbf{Response bias} (or Activity bias): Reviews on Amazon, Twitter
  tweets, Facebook comments etc. are more often written by people with
  very strong (often negative) opinions, and usually do not reflect the
  opinions of the population as a whole
\end{itemize}
\end{frame}

\begin{frame}[fragile]{Hidden Features}
\phantomsection\label{hidden-features}
\begin{itemize}
\tightlist
\item
  Also known as \emph{confounder variables} (or \emph{mediator
  variables} in psychology)
\item
  These are variables that influence your target outcome, but that you
  didn't measure
\item
  Simple example:

  \begin{itemize}
  \tightlist
  \item
    A customer happiness survey done over 7 days concludes that
    customers are very unhappy on Fridays and Saturdays
  \item
    Hidden feature: There was a rainstorm on these two days
  \item
    The hidden feature \texttt{weather}, not the weekday, was
    responsible for the drop in happiness
  \end{itemize}
\item
  As a more realistic example, data measured over multiple facilities
  may contain a similar bias, when measurement devices have slight
  differences in those locations
\end{itemize}
\end{frame}

\begin{frame}{Pipeline}
\phantomsection\label{pipeline}

%TODO put this closer to data soures, include graphics of an example
\begin{itemize}
\tightlist
\item
  You usually need to set up a pipeline to ingest new data and refresh
  your existing data set regularly.
\item
  A pipeline can be

  \begin{itemize}
  \tightlist
  \item
    batch-based
  \item
    streaming or real-time
  \item
    a hybrid of these two
  \end{itemize}
\end{itemize}
\end{frame}

\begin{frame}{Data Versioning}
\phantomsection\label{data-versioning}
\begin{itemize}
\tightlist
\item
  Your data set may change over time
\item
  Like version control for code (e.g.~git or subversion), it can make
  sense to have version control for your data sets
\item
  \emph{Pachyderm} and \emph{DVC} (Data Version Control) are two
  versioning tools
\end{itemize}
\end{frame}

\begin{frame}{Data Governance and Security}
\phantomsection\label{data-governance-and-security}
%TODO legal aspects, privacy? How to anonymize data?
\begin{itemize}
\tightlist
\item
  Good data governance protects the integrity of your business's data
\item
  Ideally, you assign responsibilities and accountabilities to ensure
  your data stays \textbf{secure, accurate, and usable}
\item
  Establish internal rules for data use
\item
  Ideally, all your data is collected into one location

  \begin{itemize}
  \tightlist
  \item
    This makes access control much easier
  \end{itemize}
\end{itemize}
\end{frame}

\begin{frame}{Data Governance - Integrity}
\phantomsection\label{data-governance---integrity}
\begin{itemize}
\tightlist
\item
  Corrupt data severely decreases its value
\item
  When your data gets updated, make sure that it stays intact
\item
  Wrong updates can add duplicates or contradictory data points, which
  invalidate your model
\item
  You can e.g.~run tests and sanity checks each time your data is
  updated
\end{itemize}
\end{frame}

\begin{frame}{Data Governance - Access}
\phantomsection\label{data-governance---access}
\begin{itemize}
\tightlist
\item
  You have to be able to access your data at all times
\item
  Ensure your data pipeline is robust and stable and it can't easily
  break
\item
  Minimize outage risks by creating a simple and robust infrastructure
\end{itemize}
\end{frame}

\begin{frame}{Wrangling, Exploration and Cleaning}
\phantomsection\label{wrangling-exploration-and-cleaning}
\begin{itemize}
\tightlist
\item
  Once you successfully ingested the necessary data, you can begin to
  prepare it for modeling
\item
  This preparation often takes up most of the time in a ML project
\item
  The necessary steps include exploratory data analysis (EDA), data
  wrangling (e.g.~joins) and data cleaning
\end{itemize}
\end{frame}

\begin{frame}{Data Labeling}
\phantomsection\label{data-labeling}
\begin{itemize}
\tightlist
\item
  Sometimes the data you obtain is not yet labeled
\item
  For example, downloading random images from the internet that you want
  to classify or tag for a computer vision project
\item
  In these cases, you need humans to assign labels to your data points
\item
  It's tempting to use cheap and/or crowdsourced workers, but be careful
  if you need domain knowledge for correct labeling!
\end{itemize}
\end{frame}

\section{Modeling}\label{modeling-1}

%TODO include overview over modelling steps (graphics?)
%TODO give overview over train-test-tune process

\begin{frame}{Feature Engineering}
\phantomsection\label{feature-engineering}
\begin{itemize}
\tightlist
\item
  Feature Engineering involves the selection, aggregation and
  transformation of your raw data to derive meaningful features
\item
  Successful feature engineering combines domain expertise with the
  insights you gained during your data exploration
\item
  Don't overdo it: Including too many useless features introduces
  unnecessary noise into your model\\
  Remember: You have to generate all features for new data.
\end{itemize}
\end{frame}

\begin{frame}{Model Training}
\phantomsection\label{model-training}
Consider your model training as a scientific experiment:

\begin{itemize}
\tightlist
\item
  Like in a lab experiment, you want to precisely keep track of all
  input parameters

  \begin{itemize}
  \tightlist
  \item
    Data version
  \item
    ML algorithm and hyperparameters
  \item
    Code version (e.g.~a git tag ID)
  \end{itemize}
\end{itemize}
\end{frame}

\begin{frame}{Model Training}
\phantomsection\label{model-training-1}
It's usually best to proceed in an iterative manner:

\begin{itemize}
\tightlist
\item
  \emph{Start simple}: Quickly create a baseline model
\item
  Then increase the model complexity and verify that each iteration
  improves your target measure
\end{itemize}
\end{frame}

\begin{frame}{Model Configuration}
\phantomsection\label{model-configuration}
\begin{itemize}
\tightlist
\item
  How do you store a configuration setup you used for a specific
  experiment (model)?\\
  \(\rightarrow\) And how do you store \emph{small} changes to that
  setup?
\item
  This configuration should be stored and clearly readable (i.e.~in
  plain text or json)
\item
  Good tools can help, e.g.~\href{https://hydra.cc}{hydra}
\end{itemize}
\end{frame}

\begin{frame}{Model Cards}
\phantomsection\label{model-cards}
\begin{itemize}
\tightlist
\item
  Similar to the Datasheets discussed earlier, you ideally want to
  document everything you can about your model
\item
  Create a \textbf{Model Card} that contains

  \begin{itemize}
  \tightlist
  \item
    performance characteristics
  \item
    the intended use cases
  \item
    usages that you don't recommend or support
  \end{itemize}
\item
  See also: \href{https://arxiv.org/pdf/1810.03993.pdf}{Model Cards for
  Model Reporting}
\end{itemize}
\end{frame}

\begin{frame}{Model Evaluation}
\phantomsection\label{model-evaluation}
%TODO how to split in train, test
\begin{itemize}
\tightlist
\item
  To select the best model, you evaluate your experiments based on the
  metrics you predefined in the business understanding phase
\item
  Besides this performance scoring, you may want to evaluate your
  model's robustness: How sensitive is it to noise or other pertubations
  and adverserial attacks?

  \begin{itemize}
  \tightlist
  \item
    You can add random noise to your dataset and investigate the impact
    on your model's parameters and/or performance
  \end{itemize}
\end{itemize}
\end{frame}

\begin{frame}{Interpretation}
\phantomsection\label{interpretation}
\begin{itemize}
\tightlist
\item
  To understand your model better, you can use methods from
  Interpretable Machine Learning (IML)

  \begin{itemize}
  \tightlist
  \item
    E.g. compute the \emph{feature importances} to see how much each of
    your features contributes to a model's prediction
  \end{itemize}
\item
  IML is also useful to uncover algorithmic bias, and prevent legal or
  ethical trouble

  \begin{itemize}
  \tightlist
  \item
    E.g., in October 2018
    \href{https://www.theguardian.com/technology/2018/oct/10/amazon-hiring-ai-gender-bias-recruiting-engine}{world
    headlines reported} about an Amazon AI recruiting tool that favored
    men. Amazon's model was trained on biased data that were skewed
    towards male candidates. It has built rules that penalized resumes
    that included the word ``women's''.
  \end{itemize}
\item
  Algorithmic bias \url{https://en.wikipedia.org/wiki/Algorithmic_bias}
\end{itemize}
\end{frame}

\section{Model Governance}\label{model-governance-1}

\begin{frame}{Model Governance}
\phantomsection\label{model-governance-2}
\begin{quote}
Developing and deploying ML systems is relatively fast and cheap, but
maintaining them over time is difficult and expensive.
~\href{https://papers.nips.cc/paper/5656-hidden-technical-debt-in-machine-learning-systems.pdf}{Hidden
Technical Debt in ML Systems}
\end{quote}
\end{frame}

\begin{frame}{Model Governance}
\phantomsection\label{model-governance-3}
\begin{itemize}
\tightlist
\item
  Model Governance includes the \textbf{people, processes and
  technologies} needed to manage and protect your ML model.
\item
  Two extremes have to be avoided:

  \begin{itemize}
  \tightlist
  \item
    \textbf{Repression} by too many standards, rules, and controlling
  \item
    \textbf{Chaos} by too much speed, freedom, creativity and change
  \end{itemize}
\item
  A good Governance ensures your application can be used for a long
  time, and can easily be maintained and adapted in case some
  requirements change
\end{itemize}
\end{frame}

\begin{frame}{Deployment}
\phantomsection\label{deployment}
%TODO include graphics on deployment strategies
\begin{itemize}
\tightlist
\item
  To make your model \emph{useful} and generate business value, you have
  to \emph{deploy} it, i.e.~make it accessible to its users

  \begin{itemize}
  \tightlist
  \item
    Users can be humans or automated systems (other programs)
  \end{itemize}
\item
  This is where ML skills fade into the background, and software
  engineering and \emph{DevOps} skills become important
\item
  You can deploy your model in different ways:

  \begin{itemize}
  \tightlist
  \item
    \textbf{One-off}: If you only need your model trained once (or very
    rarely), train it by hand and push it to production
  \item
    \textbf{Batch}: With batch training, you constantly (e.g.~weekly)
    update your model with a new batch of data (e.g.~this week's new
    contracts)
  \item
    \textbf{Real-time}: Some models are very time-critical and need to
    be updated continously on a stream of data.
  \end{itemize}
\end{itemize}
\end{frame}

\begin{frame}{Model Versioning}
\phantomsection\label{model-versioning}
\begin{itemize}
\tightlist
\item
  If you retrain and update your model frequently, it makes sense to
  keep a model registry
\item
  This is similar to ``classic'' version control for code
\item
  \href{https://mlflow.org/}{MLflow} is an open source platform for ML
  lifecycle management that includes a model registry
\end{itemize}
\end{frame}

\begin{frame}{Monitoring}
\phantomsection\label{monitoring}
\begin{itemize}
\tightlist
\item
  Our job is not finished when a model gets deployed
\item
  We need to \emph{monitor} the application to detect things like

  \begin{itemize}
  \tightlist
  \item
    problems/bugs in the pipeline
  \item
    shifts in the data distribution
  \item
    changes in the meaning of the data
  \end{itemize}
\end{itemize}
\end{frame}

\begin{frame}{Monitoring - Data drift}
\phantomsection\label{monitoring---data-drift}
\begin{itemize}
\tightlist
\item
  \textbf{Data drift} happens when your data evolves over time

  \begin{itemize}
  \tightlist
  \item
    Example: You train an image classifier on animal pictures taken in
    summer, and then winter comes and images start containing snow.
  \end{itemize}
\item
  Your model can exhibit undefined behavior in these regions

  \begin{itemize}
  \tightlist
  \item
    A linear model just extrapolates, which can make sense or not
    depending on the data
  \item
    A random forest often just ``randomly'' guesses a label
  \end{itemize}
\item
  Another possibility is that your \emph{target} evolves, and now
  includes new categories or new value ranges in case of a continuous
  target
\item
  Remedy: Enough of the new data needs to be labeled to introduce the
  model to the new classes, and the model needs to be retrained
\end{itemize}
\end{frame}

\begin{frame}{Monitoring - Concept drift}
\phantomsection\label{monitoring---concept-drift}
\begin{itemize}
\tightlist
\item
  \textbf{Concept drift} means that our interpretation of the data
  changes

  \begin{itemize}
  \tightlist
  \item
    Example: A classification problem into ``high blood pressure'', and
    the cutoff for ``high blood pressure'' changes from 140mmHg to
    135mmHg.
  \end{itemize}
\item
  The general distribution of the data may even stay the same
\item
  This means that your \emph{existing} data may become invalid
\item
  Remedy: Your affected old data needs to be relabeled and the model
  retrained
\end{itemize}
\end{frame}

\begin{frame}{User Feedback and Acceptance}
\phantomsection\label{user-feedback-and-acceptance}
\begin{itemize}
\tightlist
\item
  Two main tasks are necessary for this final step:

  \begin{itemize}
  \tightlist
  \item
    \textbf{System Validation}: Ask for customer feedback if the model
    and data pipeline meet their needs
  \item
    \textbf{Project handover}: Transfer the project to the team that
    will run your system in production
  \end{itemize}
\item
  You should generate a final report with all details that are useful
  for learning about how to operate the system.
\end{itemize}
\end{frame}

\begin{frame}{Feedback Loops / Self Reinforcing Behavior}
\phantomsection\label{feedback-loops-self-reinforcing-behavior}
%TODO include graphics for data flow here?
\begin{itemize}
\tightlist
\item
  A deployed ML model often influences its own behavior if it updates
  over time. This can lead to unexpected behavior
\item
  \textbf{Direct feedback loops}: A model may influence the selection of
  \emph{its own} future training data.

  \begin{itemize}
  \tightlist
  \item
    Example: Netflix influences which movies it recommends to certain
    viewers, and it will then use (only) these movies as future training
    data
  \end{itemize}
\item
  \textbf{Hidden feedback loops}: Here, \emph{two systems} influence
  each other \emph{indirectly} through the world.

  \begin{itemize}
  \tightlist
  \item
    Example: Flash crashes that occur because many trading algorithms
    sell at the same time and augment each others fear
  \end{itemize}
\end{itemize}
\end{frame}

\section{Best Practices for Competition Success}
%TODO graphics with overview over steps
\begin{frame}{Step 1: Understand the Problem Statement}
\begin{itemize}
  \item \textbf{Clarify Objectives \& Business Impact}: Determine whether it’s a classification or regression task, the domain context, and what “success” looks like.
  \item \textbf{Target Variable}: Understand how the target is defined and distributed.
  \item \textbf{Constraints \& Permissions}: Check for external data limits, privacy, or regulatory issues.
  \item \textbf{Action Points}:
    \begin{itemize}
      \item Annotate key requirements and assumptions.
      \item Research the domain for feature ideas.
      \item Restate the problem in your own words to ensure clarity.
    \end{itemize}
\end{itemize}
\end{frame}

\begin{frame}{Step 2: Study the Evaluation Metric}
\begin{itemize}
  \item \textbf{Metric Matters}: Each competition or project has a specific metric (e.g., F1-score, AUC-ROC, RMSE), which guides how you tune and compare models.
  \item \textbf{Classification vs. Regression Metrics}:
    \begin{itemize}
      \item Accuracy, F1-score, AUC-ROC, and Log Loss for classification.
      \item RMSE, MAE, R\textsuperscript{2} for regression.
    \end{itemize}
  \item \textbf{Optimization Direction}: Decide if you must maximize or minimize the metric; some metrics (e.g., AUC-ROC) may need a proxy loss for training.
  \item \textbf{Action Points}:
    \begin{itemize}
      \item Simulate prediction changes to see impact on the metric.
      \item If allowed, implement custom losses aligned to the competition metric.
    \end{itemize}
\end{itemize}
\end{frame}

\begin{frame}{Step 3: Perform an Initial Data Analysis (EDA)}
\begin{itemize}
  \item \textbf{Data Profiling}: Summarize shape, basic statistics, and distribution of features and target.
  \item \textbf{Quality Checks}: Identify missing values, outliers, or suspicious patterns; confirm no hidden data leakage.
  \item \textbf{Train-Test Shift}: Check if training and test sets differ significantly in their feature distributions.
  \item \textbf{Action Points}:
    \begin{itemize}
      \item Use tools like Pandas Profiling or Sweetviz for automated EDA.
      \item Document anomalies or domain-specific oddities.
      \item Examine correlations, scatter plots, and histograms.
    \end{itemize}
\end{itemize}
\end{frame}

\begin{frame}{Step 4: Develop a Baseline Model}
\begin{itemize}
  \item \textbf{Why a Baseline?} A simple model (e.g., Logistic Regression, Decision Tree) gives you a reference performance.
  \item \textbf{Pipeline Check}: Ensures data loading, preprocessing, and validation splits are correct.
  \item \textbf{Diagnostics}: Analyze misclassifications and residuals to spot improvement areas.
  \item \textbf{Action Points}:
    \begin{itemize}
      \item Start with a naive or default-parameter model.
      \item Use cross-validation to assess reliability.
      \item Record baseline metrics to quantify future gains.
    \end{itemize}
\end{itemize}
\end{frame}

\begin{frame}{Step 5: Set Up a Robust Validation Strategy}
\begin{itemize}
  \item \textbf{Importance of Validation}: Avoid overfitting to the training set or a single random split.
  \item \textbf{Common Methods}:
    \begin{itemize}
      \item K-Fold (balanced data).
      \item Stratified K-Fold (imbalanced classes).
      \item Time-Based or Group Splits (time series, repeated measures).
    \end{itemize}
  \item \textbf{Action Points}:
    \begin{itemize}
      \item Match validation approach to the final test scenario.
      \item Keep a holdout set for final checks.
      \item Use consistent folds across experiments to compare fairly.
    \end{itemize}
\end{itemize}
\end{frame}

\begin{frame}{Step 6: Feature Engineering and Data Preprocessing}
\begin{itemize}
  \item \textbf{Missing Data}: Mean/median imputation, advanced methods, or flags for “missingness.”
  \item \textbf{Encoding Categorical Variables}: One-hot, target, or frequency encoding (watch for leakage).
  \item \textbf{Scaling \& Transformations}: Normalize or log-transform skewed features, especially for sensitive models (SVMs, neural nets).
  \item \textbf{New Features}: Time-based, domain knowledge, interactions (e.g., multiplications or ratios).
  \item \textbf{Action Points}:
    \begin{itemize}
      \item Add features in small batches and track validation changes.
      \item Use scikit-learn Pipelines to avoid leakage.
      \item Consider dimensionality reduction if feature space is large.
    \end{itemize}
\end{itemize}
\end{frame}

\begin{frame}{Step 7: Select and Tune Models Strategically}
\begin{itemize}
  \item \textbf{Model Choices}:
    \begin{itemize}
      \item Tree-based ensembles (XGBoost, LightGBM, CatBoost) for tabular data.
      \item Neural networks for large-scale or unstructured data.
      \item Simple linear/logistic models can still compete with good feature engineering.
    \end{itemize}
  \item \textbf{Hyperparameter Tuning}:
    \begin{itemize}
      \item Grid Search, Random Search, Bayesian Optimization.
      \item Early stopping and regularization (L1, L2) to avoid overfitting.
    \end{itemize}
  \item \textbf{Action Points}:
    \begin{itemize}
      \item Track experiments (MLflow, Weights \& Biases, or spreadsheets).
      \item Start with defaults and refine gradually.
      \item Use interpretability tools (SHAP, permutation importance) for feature insights.
    \end{itemize}
\end{itemize}
\end{frame}

\begin{frame}{Step 8: Track Experiments and Iterate}
\begin{itemize}
  \item \textbf{Experiment Logging}: Keep detailed records of data versions, hyperparams, code revisions.
  \item \textbf{Error Analysis}: Study misclassified samples or large residuals to guide feature tweaks.
  \item \textbf{Iterative Improvements}:
    \begin{itemize}
      \item Refine features, adjust validation, revisit transformations based on your findings.
      \item Maintain a “changelog” for each iteration.
    \end{itemize}
  \item \textbf{Action Points}:
    \begin{itemize}
      \item Automate repeated tasks (submissions, data prep).
      \item Prune ineffective ideas swiftly.
      \item Ensure reproducibility (version control, environment snapshots).
    \end{itemize}
\end{itemize}
\end{frame}

\begin{frame}{Step 9: Understand the Leaderboard and Avoid Overfitting}
\begin{itemize}
  \item \textbf{Public vs. Private Leaderboard}: Public boards use only part of the test set—overfitting leads to dramatic rank changes later.
  \item \textbf{Submission Strategy}:
    \begin{itemize}
      \item Rely mainly on internal validation.
      \item Keep a private holdout to confirm final performance.
    \end{itemize}
  \item \textbf{Leaderboard Shakeups}: Watch for big changes when private scores are revealed; it often indicates overfitting to the public set.
  \item \textbf{Action Points}:
    \begin{itemize}
      \item Limit minor tweaks solely to improve a public rank.
      \item If a jump seems suspicious, re-check data handling for leakage.
    \end{itemize}
\end{itemize}
\end{frame}

\begin{frame}{Step 10: Collaborate and Learn from Others}
\begin{itemize}
  \item \textbf{Study Top Solutions}: Many winning approaches are documented in blog posts or Kaggle discussions.
  \item \textbf{Teamwork}:
    \begin{itemize}
      \item Collaboration brings diverse perspectives and model ensembling.
      \item Share code and features for synergy.
    \end{itemize}
  \item \textbf{Community \& Networking}:
    \begin{itemize}
      \item Discussion forums reveal dataset quirks, best practices.
      \item Local meetups or online communities keep you updated.
    \end{itemize}
  \item \textbf{Action Points}:
    \begin{itemize}
      \item Document and share your findings.
      \item Adapt others’ ideas carefully—validate them on your own splits.
      \item Engage with peers for new tools and techniques.
    \end{itemize}
\end{itemize}
\end{frame}

\begin{frame}{Final Advice}
\begin{itemize}
  \item \textbf{Start Simple, Then Iterate}: Don’t jump into complex models before a solid baseline.
  \item \textbf{Prioritize Validation}: A robust validation strategy prevents nasty surprises.
  \item \textbf{Log Everything}: Document all transformations, parameter settings, and model versions.
  \item \textbf{Aim for Generalizable Solutions}: Don’t just chase leaderboard scores—focus on reproducibility and reliability.
\end{itemize}
\end{frame}



\end{document}
